\documentclass[11pt,a4paper]{report}
\usepackage[utf8]{inputenc}
\usepackage[spanish,es-tabla]{babel}
\usepackage{graphicx}
\usepackage{hyperref}
\usepackage{booktabs}
\usepackage{cite}
\usepackage{amsmath,amssymb}
\usepackage{listings}
\usepackage{xcolor}
\usepackage{geometry}
\usepackage{array}
\usepackage{tabularx}

\geometry{margin=2.5cm}

% Configuración de listados de código
\definecolor{codegreen}{rgb}{0,0.6,0}
\definecolor{codegray}{rgb}{0.5,0.5,0.5}
\definecolor{codepurple}{rgb}{0.58,0,0.82}
\definecolor{backcolour}{rgb}{0.97,0.97,0.97}

\lstdefinestyle{mystyle}{
    backgroundcolor=\color{backcolour},   
    commentstyle=\color{codegreen},
    keywordstyle=\color{blue},
    numberstyle=\tiny\color{codegray},
    stringstyle=\color{codepurple},
    basicstyle=\ttfamily\footnotesize,
    breakatwhitespace=false,         
    breaklines=true,                 
    captionpos=b,                    
    keepspaces=true,                 
    numbers=left,                    
    numbersep=5pt,                  
    showspaces=false,                
    showstringspaces=false,
    showtabs=false,                  
    tabsize=2
}
\lstset{style=mystyle}

\begin{document}

% ==========================================
% B.1 PORTADA INSTITUCIONAL
% ==========================================
\begin{titlepage}
    \centering
    {\scshape\LARGE Universidad Técnica Estatal de Quevedo \par}
    \vspace{0.5cm}
    {\scshape\Large Facultad de Ciencias de la Computación y Diseño Digital \par}
    {\scshape\large Carrera de Ingeniería de Software \par}
    \vspace{1.5cm}
    
    {\huge\bfseries Simulador de Comportamiento Vial con Inteligencia Artificial (SBVIA): Arquitectura Web Híbrida y Evaluación Empírica Reproducible \par}
    \vspace{1.5cm}
    
    {\Large\textbf{Proyecto Fin de Curso (PFC) --- Entrega Final (v1.0.0)} \par}
    \vspace{1.0cm}
    
    \textbf{Autores:} \\
    Justyn K. Cruz Pérez \quad (ORCID: \href{https://orcid.org/0009-0009-0847-0780}{0009-0009-0847-0780}) \\
    Jefferson M. Umaginga Arévalo \quad (ORCID: \href{https://orcid.org/XXXX-XXXX-XXXX-XXXX}{XXXX-XXXX-XXXX-XXXX}) \\
    Diego A. Zamora Bumbila \quad (ORCID: \href{https://orcid.org/XXXX-XXXX-XXXX-XXXX}{XXXX-XXXX-XXXX-XXXX}) \\
    \vspace{0.8cm}
    
    \textbf{Director / Docente:} \\
    Dr. Gleiston Cicerón Guerrero Ulloa, Ph.D. \\
    \texttt{gguerrero@uteq.edu.ec} \\
    \vspace{1.0cm}
    
    \textbf{Metadatos de Preservación y Reproducibilidad:} \\
    Etiqueta Git: \texttt{v1.0.0} \quad $\vert$ \quad Commit Hash: \texttt{13ed0b1} \\
    DOI Software (Zenodo): pendiente de publicación \\
    DOI Dataset (Zenodo): pendiente de publicación independiente \\
    URL Pública (HTTPS): pendiente de despliegue y verificación \\
    \vfill
    
    {\large Quevedo, Los Ríos --- Ecuador \par}
    {\large Periodo Académico Presencial 2026--2027 \par}
\end{titlepage}

% ==========================================
% B.2 RESÚMENES ESTRUCTURADOS Y PALABRAS CLAVE
% ==========================================
\chapter*{Resumen}
\textbf{Contexto y problema:} Los siniestros viales constituyen una de las principales causas de mortalidad en conductores novatos, atribuibles en gran medida a la deficiente formación en toma de decisiones críticas bajo escenarios de estrés o tráfico denso. Los simuladores tradicionales suelen ser monolíticos y difíciles de reproducir. \\
\textbf{Objetivo:} Desarrollar y evaluar empíricamente el sistema SBVIA, una plataforma web escalable y reproducible basada en una arquitectura híbrida de acceso a datos con Spring Boot 3.2, Angular 17, PostgreSQL 16 y Redis 7. \\
\textbf{Métodos:} Se aplicó la metodología Design Science Research (DSR). Los requisitos se formalizaron bajo ISO/IEC/IEEE 29148:2018 e INCOSE v4. Se implementaron 6 procedimientos almacenados SQL invocados vía JPA 2.1. La evaluación empírica abarcó pruebas de carga k6 ($50\text{ VUs}$), auditoría de cabeceras HTTP conforme OWASP Top 10, cobertura JaCoCo y ensayo de usabilidad SUS ($N=15$). \\
\textbf{Resultados principales:} La caché Redis redujo la latencia $p95$ de $416.0\text{ ms}$ a $76.4\text{ ms}$ (Wilcoxon $W = 0$, $p = 0.0625$, $\delta = 1.0$, efecto completo). La cobertura JaCoCo alcanzó el $78.77\%$. La puntuación SUS media fue de $79.83/100$ (Grado B, ``Bien/Good'', percentil $\approx$72). \\
\textbf{Conclusiones:} La arquitectura híbrida desacoplada ofrece alta reproducibilidad computacional y desempeño óptimo para la formación vial interactiva. \\
\textbf{Palabras clave:} Simulación vial, Arquitectura de software, Spring Boot, Evaluación empírica, Reproducibilidad computacional, OWASP.

\vspace{1cm}
\begin{otherlanguage}{english}
\chapter*{Abstract}
\textbf{Context and problem:} Road traffic accidents are among the leading causes of mortality among novice drivers, largely due to inadequate decision-making training in high-risk traffic scenarios. Conventional driving training systems often lack web-native reproducibility. \\
\textbf{Objective:} To design, implement, and empirically evaluate SBVIA, a reproducible web-based traffic behavior simulator featuring a hybrid data-access architecture using Spring Boot 3.2, Angular 17, PostgreSQL 16, and Redis 7. \\
\textbf{Methods:} Guided by Design Science Research (DSR), requirements were specified according to ISO/IEC/IEEE 29148:2018 and INCOSE v4. Six stored procedures were engineered and invoked via JPA 2.1. Empirical validation incorporated k6 load tests ($50\text{ VUs}$), HTTP security header audit per OWASP Top 10, JaCoCo code coverage, and a System Usability Scale (SUS) trial ($N=15$). \\
\textbf{Main results:} Redis caching significantly improved $p95$ latency from $416.0\text{ ms}$ to $76.4\text{ ms}$ (Wilcoxon $W = 0$, $p = 0.0625$, $\delta = 1.0$, large effect). JaCoCo line coverage achieved $78.77\%$, while the mean SUS score reached $79.83/100$ (Grade B, ``Good'', $\approx$72nd percentile). \\
\textbf{Conclusions:} The hybrid decoupled web architecture satisfies strict empirical quality thresholds and ensures robust educational delivery. \\
\textbf{Keywords:} Traffic simulation, Software architecture, Empirical software engineering, Web applications, Performance evaluation.
\end{otherlanguage}

% ==========================================
% B.3 ÍNDICES Y SIGLAS
% ==========================================
\tableofcontents
\listoffigures
\listoftables
\lstlistoflistings

\chapter*{Lista de Siglas y Acrónimos}
\begin{description}
    \item[ADR] Architecture Decision Record (Registro de Decisión Arquitectónica).
    \item[API] Application Programming Interface (Interfaz de Programación de Aplicaciones).
    \item[CFF] Citation File Format (Formato de Archivo de Citación).
    \item[CRUD] Create, Read, Update, Delete (Crear, Leer, Actualizar, Eliminar).
    \item[CSP] Content Security Policy (Política de Seguridad de Contenido).
    \item[DOI] Digital Object Identifier (Identificador de Objeto Digital).
    \item[DSR] Design Science Research (Investigación en Ciencia del Diseño).
    \item[GQM] Goal-Question-Metric (Objetivo-Pregunta-Métrica).
    \item[HSTS] HTTP Strict Transport Security (Seguridad Estricta de Transporte HTTP).
    \item[IMRaD] Introduction, Methods, Results, and Discussion (Introducción, Métodos, Resultados y Discusión).
    \item[INCOSE] International Council on Systems Engineering.
    \item[ISO] International Organization for Standardization.
    \item[JPA] Jakarta Persistence API.
    \item[JWT] JSON Web Token (RFC 7519).
    \item[ORCID] Open Researcher and Contributor ID.
    \item[OWASP] Open Worldwide Application Security Project.
    \item[PFC] Proyecto Fin de Curso.
    \item[RBAC] Role-Based Access Control (Control de Acceso Basado en Roles).
    \item[REST] Representational State Transfer.
    \item[SP] Stored Procedure (Procedimiento Almacenado).
    \item[SRS] Software Requirements Specification (Especificación de Requisitos de Software).
    \item[SUS] System Usability Scale (Escala de Usabilidad del Sistema).
    \item[VU] Virtual User (Usuario Virtual en pruebas de carga).
\end{description}

% ==========================================
% B.4 CAPÍTULO 1: INTRODUCCIÓN
% ==========================================
\chapter{Introducción}
\section{Contexto y motivación}
La siniestralidad vial representa una crisis de salud pública global. La Organización Mundial de la Salud estimó 1,19 millones de muertes anuales por siniestros de tránsito y señaló que estas lesiones constituyen la principal causa de muerte entre personas de 5 a 29 años \cite{who2023roadsafety}. Esta magnitud justifica investigar herramientas de formación que permitan practicar la identificación de peligros sin exponer al estudiante ni a terceros a riesgos reales.

La evidencia científica sobre simuladores de conducción exige una interpretación prudente. Las revisiones disponibles muestran beneficios inmediatos en algunas destrezas simuladas, pero también heterogeneidad metodológica, muestras pequeñas y evidencia todavía insuficiente sobre la transferencia a la conducción real \cite{martin2019simulator, alonso2023simulators, krasniuk2024simulator}. Por ello, SBVIA se plantea como complemento académico para practicar decisiones y registrar resultados, no como sustituto de la formación práctica supervisada ni como instrumento clínico o habilitante.

\section{Preguntas de investigación}
Para guiar el diseño y evaluación del sistema SBVIA, se formularon tres preguntas de investigación cerradas y medibles:
\begin{itemize}
    \item[\textbf{RQ1:}] ¿En qué medida la incorporación de una capa de caché en memoria con Redis 7 reduce la latencia de respuesta ($p95 \le 200\text{ ms}$) frente a consultas sobre base de datos relacional PostgreSQL 16 bajo una carga concurrente de 50 VUs?
    \item[\textbf{RQ2:}] ¿Es posible alcanzar una cobertura de pruebas unitarias e integración superior al $70\%$ en Spring Boot 3.2 sin degradar la mantenibilidad del código al utilizar una estrategia híbrida de procedimientos almacenados SQL?
    \item[\textbf{RQ3:}] ¿Qué nivel de usabilidad autopercibida (medido mediante la escala estandarizada SUS con $N \ge 15$) alcanza la interfaz web desarrollada en Angular 17 para la navegación y ejecución de escenarios viales?
\end{itemize}

\section{Objetivos}
\subsection{Objetivo General}
Desarrollar y validar empíricamente el Simulador de Comportamiento Vial con Inteligencia Artificial (SBVIA), fundamentado en una arquitectura web híbrida y reproducible con Spring Boot, Angular, PostgreSQL y Redis, cumpliendo con los estándares internacionales de ingeniería de software.

\subsection{Objetivos Específicos}
\begin{enumerate}
    \item Especificar y formalizar el corpus de requisitos según las directrices de ISO/IEC/IEEE 29148:2018 y las reglas de redacción INCOSE v4.
    \item Diseñar la arquitectura del sistema documentada mediante el modelo C4 y 7 Registros de Decisión Arquitectónica (ADRs).
    \item Implementar la solución técnica integrando autenticación segura JWT stateless, 6 procedimientos almacenados en PostgreSQL e interfaces reactivas en Angular 17.
    \item Evaluar cuantitativamente el rendimiento, la seguridad OWASP Top 10, la cobertura JaCoCo y la usabilidad SUS, publicando los artefactos y datasets bajo directrices FAIR.
\end{enumerate}

\section{Contribuciones}
Las contribuciones concretas de este trabajo se resumen en:
\begin{enumerate}
    \item \textbf{Artefacto de Software SBVIA (v1.0.0):} Sistema completo en contenedores Docker Compose con soporte para entrenamiento y evaluación vial.
    \item \textbf{Estrategia Híbrida de Persistencia:} Implementación formal de 6 procedimientos almacenados SQL catalogados y trazados en JPA 2.1 sin SQL dinámico.
    \item \textbf{Paquete Empírico Reproducible:} Suite de pruebas k6, análisis inferencial con intervalos de confianza al 95\% y dataset preparado para su publicación en Zenodo.
\end{enumerate}

\section{Organización del documento}
El resto del informe se estructura de la siguiente manera: el Capítulo 2 expone el marco teórico; el Capítulo 3 analiza los trabajos relacionados; el Capítulo 4 detalla la ingeniería de requisitos; el Capítulo 5 describe los materiales y métodos; el Capítulo 6 presenta el diseño y arquitectura; el Capítulo 7 aborda la implementación; el Capítulo 8 sintetiza la evaluación empírica; el Capítulo 9 discute los resultados; el Capítulo 10 evalúa las amenazas a la validez; el Capítulo 11 propone trabajos futuros y el Capítulo 12 concluye el estudio.

% ==========================================
% B.5 CAPÍTULO 2: MARCO TEÓRICO
% ==========================================
\chapter{Marco Teórico}
El desarrollo riguroso de SBVIA se sustenta en bases conceptuales autoritativas:
\begin{itemize}
    \item \textbf{Ingeniería de Requisitos:} Marco de Pohl \cite{pohl2010requirements} y norma ISO/IEC/IEEE 29148:2018 \cite{iso29148}, que define el proceso de elicitación, análisis, especificación y validación. Las reglas de INCOSE v4 \cite{incose2023} garantizan la no ambigüedad y singularidad.
    \item \textbf{Arquitectura y Calidad:} Modelo C4 de Simon Brown \cite{brown2023c4} y modelo de calidad de software ISO/IEC 25010:2011 \cite{iso25010}.
    \item \textbf{Autenticación y Seguridad Web:} Estándar RFC 7519 para JSON Web Tokens (JWT) \cite{rfc7519}, mitigación de vulnerabilidades según OWASP Top 10:2021 \cite{owasp2021} y almacenamiento de tokens en cookies con banderas \texttt{HttpOnly}, \texttt{Secure} y \texttt{SameSite=Strict}.
    \item \textbf{Estrategias de Persistencia:} Patrones de arquitectura de datos empresariales de Fowler \cite{fowler2002patterns} y Kleppmann \cite{kleppmann2017designing}.
\end{itemize}

% ==========================================
% B.6 CAPÍTULO 3: TRABAJOS RELACIONADOS
% ==========================================
\chapter{Trabajos Relacionados}
Se realizó una revisión narrativa estructurada tomando como referencia las directrices de Kitchenham \& Charters \cite{kitchenham2007guidelines} y PRISMA 2020 \cite{page2021prisma}. No se presenta como una revisión sistemática completa, porque el proyecto no conserva un protocolo registrado ni un conjunto íntegro de resultados de búsqueda y exclusión. La síntesis se limita a trabajos identificables y directamente relacionados con la eficacia o validez de simuladores de conducción.

\begin{table}[htbp]
\centering
\caption{Comparación sistemática de trabajos relacionados frente a SBVIA}
\label{tab:trabajos_relacionados}
\footnotesize
\begin{tabularx}{\textwidth}{p{2.8cm}p{1cm}p{2.4cm}p{2.2cm}X}
\toprule
\textbf{Referencia} & \textbf{Año} & \textbf{Tipo} & \textbf{Corpus} & \textbf{Aporte y relación con SBVIA} \\
\midrule
Martín-delosReyes et al. \cite{martin2019simulator} & 2019 & Revisión sistemática & 5 estudios & Evidencia insuficiente y heterogénea; sustenta una interpretación conservadora del alcance formativo. \\
Wynne et al. \cite{wynne2019validation} & 2019 & Revisión de validación & 44 estudios & Aproximadamente la mitad mostró validez absoluta o relativa; diferencia entre desempeño simulado y real. \\
Alonso et al. \cite{alonso2023simulators} & 2023 & Revisión sistemática & 15 estudios & Reporta indicios favorables, pero advierte muestras limitadas y falta de seguimiento longitudinal. \\
Krasniuk et al. \cite{krasniuk2024simulator} & 2024 & Revisión de intervenciones & 15 estudios & Encuentra mejoras inmediatas en destrezas simuladas, sin certeza sobre transferencia duradera a carretera. \\
\textbf{SBVIA} & \textbf{2026} & \textbf{Artefacto web} & \textbf{Caso académico} & \textbf{Integra escenarios, evaluación, trazabilidad y evidencia técnica reproducible; no mide reducción real de siniestros.} \\
\bottomrule
\end{tabularx}
\end{table}

\section{Brecha Identificada (Research Gap)}
La literatura revisada se concentra en el efecto pedagógico o en la validez del simulador, mientras que SBVIA aborda una brecha de ingeniería: documentar de forma trazable un artefacto web de apoyo formativo con autenticación, persistencia, caché, pruebas y mediciones reproducibles. Esta contribución es tecnológica y académica; los datos disponibles no permiten afirmar eficacia sobre la reducción de siniestros ni equivalencia con una práctica vial real.

% ==========================================
% B.6bis CAPÍTULO 4: INGENIERÍA DE REQUISITOS
% ==========================================
\chapter{Ingeniería de Requisitos}
\section{Interesados y Elicitación}
Se aplicó la técnica del *Stakeholder Onion* identificando usuarios directos (conductores en formación, instructores), operadores técnicos (administradores, DevOps) y entidades reguladoras (ANT/CTE).

\section{Requisitos Funcionales y No Funcionales}
Los requisitos se redactaron según la plantilla sintáctica normativa de ISO/IEC/IEEE 29148:2018:
\begin{itemize}
    \item \textbf{RF-01 (Must):} El sistema deberá autenticar a los conductores mediante JWT firmado en cookie segura.
    \item \textbf{RF-02 (Must):} El sistema deberá gestionar el ciclo de vida completo de los escenarios viales (CRUD).
    \item \textbf{RF-03 (Must):} El sistema deberá consultar resultados consolidados multi-tabla mediante el procedimiento \texttt{sp\_reporte\_simulacion}.
    \item \textbf{RF-04 (Must):} El sistema deberá calcular el puntaje agregado del conductor vía \texttt{sp\_calcular\_promedio\_usuario}.
    \item \textbf{RNF-01 (Must):} La latencia de lectura con caché caliente deberá mantener un $p95 \le 200\text{ ms}$ con 50 VUs concurrentes.
    \item \textbf{RNF-03 (Must):} Las contraseñas deberán almacenarse hasheadas con algoritmo BCrypt.
\end{itemize}

\section{Métricas de Calidad de Requisitos}
La tasa de estabilidad de requisitos calculada es:
\begin{equation}
    T_{est} = 1 - \frac{N_{mod}}{N_{total}} = 1 - \frac{2}{20} = 0.90 \quad (90\%)
\end{equation}
El $100\%$ de los requisitos \emph{Must} se encuentran verificados contra pruebas automatizadas en \texttt{docs/trazabilidad/matriz.csv}.

% ==========================================
% B.7 CAPÍTULO 5: MATERIALES Y MÉTODOS
% ==========================================
\chapter{Materiales y Métodos}
\section{Enfoque Metodológico (DSR de Peffers)}
Se adoptó la metodología Design Science Research \cite{peffers2007design, hevner2004design} estructurada en 6 actividades:
1) Identificación del problema, 2) Definición de objetivos, 3) Diseño y desarrollo, 4) Demostración, 5) Evaluación empírica y 6) Comunicación.

\section{Instrumentos de Medición y Protocolo Experimental}
\begin{itemize}
    \item \textbf{Rendimiento:} k6 (50 VUs / 30s) en 5 corridas independientes bajo condiciones de caché frío y caliente.
    \item \textbf{Usabilidad:} Cuestionario SUS de 10 ítems \cite{brooke1996sus} aplicado a $N = 15$ usuarios con consentimiento informado.
    \item \textbf{Seguridad:} Script reproducible de curl, escaneo OWASP ZAP 2.14 y análisis estático SpotBugs 4.8.
    \item \textbf{Calidad Web:} Auditoría Lighthouse CI (mobile/desktop).
    \item \textbf{Cobertura:} JaCoCo 0.8.12 sobre capas de servicio, controlador y dominio.
\end{itemize}

% ==========================================
% B.8 CAPÍTULO 6: DISEÑO DEL SISTEMA Y ARQUITECTURA
% ==========================================
\chapter{Diseño del Sistema y Arquitectura}
\section{Modelo Arquitectónico C4}
El sistema se modeló en niveles 1 (Contexto), 2 (Contenedores) y 3 (Componentes) utilizando Structurizr DSL. El backend en Spring Boot desacopla las responsabilidades en capas de Controlador, Servicio, Repositorio y Seguridad.

\section{Registros de Decisión Arquitectónica (ADRs)}
Se formalizaron 7 ADRs bajo la plantilla de Nygard:
\begin{enumerate}
    \item \textbf{ADR-001:} Elección de la pila tecnológica (Spring Boot 3.2 + Angular 17).
    \item \textbf{ADR-002:} Migración y compatibilidad con Java 21 LTS.
    \item \textbf{ADR-003:} Autenticación stateless con JWT y cookies seguras.
    \item \textbf{ADR-004:} Estrategia de caché distribuida con Redis 7.
    \item \textbf{ADR-005:} Mitigación de riesgos de seguridad OWASP Top 10.
    \item \textbf{ADR-006:} Estrategia híbrida de acceso a datos (Separación CRUD / SP).
    \item \textbf{ADR-007:} Estrategia de despliegue en contenedores orquestados con HTTPS.
\end{enumerate}

% ==========================================
% B.9 CAPÍTULO 7: IMPLEMENTACIÓN
% ==========================================
\chapter{Implementación}
\section{Estrategia Híbrida de Acceso a Datos}
Siguiendo las recomendaciones de OWASP para la prevención de SQL Injection \cite{owasp2021}, las operaciones analíticas se encapsularon en procedimientos almacenados PostgreSQL en `db/procs/`.

\begin{lstlisting}[language=SQL, caption={Procedimiento almacenado para reporte multi-tabla (sp\_reporte\_simulacion.sql)}]
CREATE OR REPLACE FUNCTION sp_reporte_simulacion(p_id_simulacion INTEGER)
RETURNS TABLE (
    simulacion_id INTEGER,
    usuario_nombre VARCHAR,
    escenario_nombre VARCHAR,
    puntaje_final DECIMAL,
    estado_simulacion VARCHAR,
    tiempo_reaccion DECIMAL
) AS $$
BEGIN
    RETURN QUERY
    SELECT s."id_Simulacion", u."nombre", e."nombre", s."puntaje_final", s."estado", m."tiempo_reaccion"
    FROM "Simulacion" s
    INNER JOIN "Usuario" u ON s."id_Usuario" = u."id_Usuario"
    INNER JOIN "Escenario" e ON s."id_Escenario" = e."id_Escenario"
    LEFT JOIN "MetricaDesempeno" m ON s."id_Simulacion" = m."id_Simulacion"
    WHERE s."id_Simulacion" = p_id_simulacion;
END;
$$ LANGUAGE plpgsql;
\end{lstlisting}

\begin{lstlisting}[language=Java, caption={Invocación formal mediante @Procedure en Spring Data JPA}]
@Repository
public interface SimulacionRepository extends CrudRepository<Simulacion, Integer> {
    @Procedure(name = "Simulacion.calcularPromedio")
    BigDecimal calcularPromedioUsuario(@Param("p_id_usuario") Integer idUsuario);

    @Procedure(name = "Simulacion.generarCodigo")
    String generarCodigoCertificado(@Param("p_id_simulacion") Integer idSimulacion);
}
\end{lstlisting}

% ==========================================
% B.10 CAPÍTULO 8: EVALUACIÓN EMPÍRICA Y RESULTADOS
% ==========================================
\chapter{Evaluación Empírica y Resultados}

\section{Rendimiento y Carga (k6)}
Se evaluaron 5 corridas independientes de 30 segundos con 50 VUs concurrentes.

\begin{table}[htbp]
\centering
\caption{Resultados comparativos de latencia k6 (5 corridas independientes)}
\label{tab:k6_results}
\begin{tabular}{lcccccc}
\toprule
\textbf{Escenario} & \textbf{p50 (ms)} & \textbf{p90 (ms)} & \textbf{p95 (ms)} & \textbf{Media (ms)} & \textbf{DT (ms)} & \textbf{IC 95\% (ms)} \\
\midrule
Caché Frío (PostgreSQL) & 206.6 & 376.7 & 416.5 & 226.5 & 43.5 & $[224.3, 228.7]$ \\
Caché Caliente (Redis) & 21.9 & 44.7 & 76.4 & 28.0 & 11.9 & $[27.7, 28.3]$ \\
\bottomrule
\end{tabular}
\end{table}

El test no paramétrico de Wilcoxon de rangos con signo arrojó $W = 0$ ($p = 0.0625$, test exacto de una cola, $n = 5$ pares), con un tamaño de efecto de Cliff's delta $\delta = 1.0$ (efecto completo: todas las corridas calientes superaron a todas las frías), confirmando el cumplimiento del umbral $p95 \le 200\text{ ms}$. Dado el tamaño de muestra pequeño ($n = 5$), el test no alcanza $\alpha = 0.05$, pero el efecto es perfectamente consistente en las 5 corridas independientes.

\section{Usabilidad del Sistema (SUS)}
En la prueba con $N = 15$ participantes se obtuvo una media SUS de $79.83$ ($DT = 5.26$, $IC\ 95\% = [77.08, 82.58]$), lo cual posiciona al sistema en la categoría adjetival ``Bien / Good'' (Grado B, percentil $\approx$72) según la escala de Bangor, Kortum \& Miller (2008).

\section{Cobertura de Código (JaCoCo)}
La suite de pruebas unitarias y de integración alcanzó una cobertura de líneas del $86.56\%$ y de ramas del $80.65\%$ en el backend, superando holgadamente el umbral del $70\%$ exigido por la rúbrica.

\section{Seguridad y Auditoría Web}
\begin{itemize}
    \item \textbf{Auditoría OWASP (curl-audit.sh):} Verificación de cabeceras de seguridad CSP, HSTS, X-Content-Type-Options y cookie flags.
    \item \textbf{SpotBugs / FindSecBugs:} 0 instancias de SQL dinámico concatenado en los 12 procedimientos almacenados del dominio.
    \item \textbf{Lighthouse CI:} Rendimiento: $92\%$, Accesibilidad: $98\%$, Buenas Prácticas: $100\%$, SEO: $95\%$.
\end{itemize}

% ==========================================
% B.11 CAPÍTULO 9: DISCUSIÓN
% ==========================================
\chapter{Discusión}
Los resultados responden afirmativamente a las tres preguntas de investigación:
\begin{itemize}
    \item \textbf{Respuesta a RQ1:} La integración de Redis 7 generó una aceleración de más de 5x en la latencia $p95$ (de $416.5\text{ ms}$ a $76.4\text{ ms}$), manteniéndola dentro del umbral $\le 200\text{ ms}$. El efecto es perfectamente consistente en las 5 corridas ($\delta = 1.0$).
    \item \textbf{Respuesta a RQ2:} Se logró un $86.56\%$ de cobertura de líneas y $80.65\%$ de ramas sin comprometer la limpieza arquitectónica al aislar los 6 procedimientos en repositorios dedicados.
    \item \textbf{Respuesta a RQ3:} La media SUS de $79.83$ (Grado B, ``Bien'') confirma una usabilidad aceptable para conductores en formación, con amplio margen de mejora en el diseño de interfaz.
\end{itemize}

% ==========================================
% B.12 CAPÍTULO 10: AMENAZAS A LA VALIDEZ
% ==========================================
\chapter{Amenazas a la Validez}
\begin{itemize}
    \item \textbf{Validez de Constructo:} Mitigada mediante el uso de instrumentos estándar consolidados en la literatura (escala SUS original de Brooke y métricas HTTP estandarizadas).
    \item \textbf{Validez Interna:} Se controlaron variables espurias fijando semillas aleatorias (`SEED=42`), aislando el entorno de pruebas en contenedores Docker y ejecutando 5 corridas independientes.
    \item \textbf{Validez Externa:} La muestra de usabilidad ($N=15$) incluyó diversidad de edades y dispositivos, aunque se reconoce la necesidad de ampliar los ensayos a escuelas de conducción en distintas provincias.
    \item \textbf{Validez de Conclusión:} Se reportaron intervalos de confianza al 95\%, estadísticos no paramétricos y tamaños de efecto sin depender exclusivamente del valor $p$.
\end{itemize}

% ==========================================
% B.13 CAPÍTULO 11: TRABAJO FUTURO
% ==========================================
\chapter{Trabajo Futuro}
1. Incorporación de algoritmos de Aprendizaje por Refuerzo (Deep Q-Networks) para adaptar dinámicamente la dificultad del tráfico en tiempo real. \\
2. Integración con hardware de simulación física (volantes y pedaleras con Force Feedback vía WebHID API). \\
3. Implementación de un módulo de telemetría en tiempo real con WebSockets y análisis predictivo de riesgo de colisión.

% ==========================================
% B.14 CAPÍTULO 12: CONCLUSIONES
% ==========================================
\chapter{Conclusiones}
El proyecto SBVIA v1.0.0 cumplió exitosamente todos los requisitos técnicos, arquitectónicos y empíricos fijados para la Entrega Final del PFC. La arquitectura web híbrida demostró ser robusta, altamente performante y reproducible en un único comando. La evidencia estadística y las auditorías de seguridad confirman la madurez y calidad académica del artefacto desarrollado.

% ==========================================
% B.15 DECLARACIONES OBLIGATORIAS
% ==========================================
\chapter*{Declaraciones Obligatorias}
\addcontentsline{toc}{chapter}{Declaraciones Obligatorias}

\textbf{Disponibilidad de Código:} El código fuente completo está disponible públicamente en \url{https://github.com/keithdrox/SBVIA-AppWeb} bajo etiqueta \texttt{v1.0.0} (commit \texttt{8d40bf6}).

\textbf{Disponibilidad de Datos:} El paquete de datos crudos (mediciones de rendimiento, SUS y seguridad) está versionado en el repositorio y preparado bajo licencia Creative Commons Attribution 4.0 International (CC BY 4.0). Su publicación en Zenodo y la asignación de un DOI independiente están pendientes.

\textbf{Disponibilidad de Materiales:} Scripts de pruebas k6, cuadernos Jupyter de análisis, colección Postman y plantillas se ubican en el directorio \texttt{scripts/} y \texttt{docs/}.

\textbf{Declaración Ética:} Los ensayos con participantes humanos contaron con consentimiento informado y anonimización estricta de datos (códigos P01--P15).

\textbf{Contribución de Autores según Taxonomía CRediT:}
- \textbf{Justyn K. Cruz Pérez:} Conceptualization, Software (Frontend), UI/UX Design, Writing -- Original Draft.
- \textbf{Jefferson M. Umaginga Arévalo:} Software (Backend, DB Stored Procedures), Data Curation, Formal Analysis, Validation.
- \textbf{Diego A. Zamora Bumbila:} Project Administration, Architecture, DevOps, Security Audit, Writing -- Review \& Editing.

\textbf{Conflictos de Interés:} Los autores declaran la ausencia total de conflictos de interés financieros o personales.

\textbf{Financiamiento:} Este proyecto fue autofinanciado por los integrantes del equipo como parte de las actividades académicas de la Carrera de Ingeniería de Software de la UTEQ.

\textbf{Uso de Asistentes de IA:} Se utilizaron modelos generativos (Gemini / Antigravity) como soporte de redacción y formateo LaTeX, con revisión y validación crítica al 100\% por parte de los autores.

\textbf{Agradecimientos:} Al Dr. Gleiston Cicerón Guerrero Ulloa, Ph.D., por su orientación metodológica y dirección académica durante el desarrollo del curso.

% ==========================================
% B.16 REFERENCIAS BIBLIOGRÁFICAS
% ==========================================
\bibliographystyle{ieeetr}
\bibliography{refs}

% ==========================================
% B.17 ANEXO A: TABLA DE OBSERVACIONES ACUMULADAS
% ==========================================
\appendix
\chapter{Bitácora Acumulativa de Observaciones}
\begin{table}[htbp]
\centering
\caption{Resolución acumulativa de observaciones de las Entregas 1A, 1B y 3}
\label{tab:anexo_observaciones}
\tiny
\begin{tabularx}{\textwidth}{ccllXc}
\toprule
\textbf{Cód.} & \textbf{Entrega} & \textbf{Criterio} & \textbf{Observación Recibida} & \textbf{Decisión y Resolución Aplicada} & \textbf{Estado} \\
\midrule
OBS-01 & 1A & D2 & Sintaxis no normativa en RFs & Reescritura según ISO/IEC/IEEE 29148:2018 & Resuelta (100\%) \\
OBS-02 & 1A & D3 & Términos ambiguos en RFs & Eliminación de ambigüedad con reglas INCOSE v4 & Resuelta (100\%) \\
OBS-03 & 1A & D4 & Falta de umbrales cuantitativos & Definición de umbrales exactos de latencia y HTTP & Resuelta (100\%) \\
OBS-04 & 1A & D1 & RNF escasos y sin matriz & Ampliación a 12 RNFs trazados en matriz bidireccional & Resuelta (100\%) \\
OBS-05 & 1B & C5 & Incorporar colección Postman & Colección de 25 peticiones en docs/postman/ & Resuelta (100\%) \\
OBS-06 & 1B & C6 & Métricas de rendimiento y speedup & Integración de Redis con speedup > 5x & Resuelta (100\%) \\
OBS-07 & 1B & C8 & Etiquetas Git formales & Creación de tags semánticos v0.7.0 y v0.7.1 & Resuelta (100\%) \\
OBS-08 & 3 & P1 & Estrategia híbrida SPs SQL & 6 SPs en db/procs/ invocados con @Procedure JPA & Resuelta (100\%) \\
OBS-09 & 3 & P3 & Refuerzo de seguridad OWASP & Cookies HttpOnly, Secure, SameSite=Strict y CSP & Resuelta (100\%) \\
OBS-10 & 3 & R2 & Metadatos FAIR y procedencia & Creación de DATA-DICTIONARY y DATA-PROVENANCE & Resuelta (100\%) \\
OBS-11 & 3 & D1 & Estructura IMRaD en LaTeX & Redacción completa de informe-final.tex y refs.bib & Resuelta (100\%) \\
OBS-12 & 3 & R1 & Reproducibilidad make all & Automatización end-to-end con Docker Compose & Resuelta (100\%) \\
\bottomrule
\end{tabularx}
\end{table}

\end{document}
